\documentclass[11pt,a4paper]{article}
\usepackage[UTF8]{ctex}
\usepackage{geometry}
\geometry{left=2.5cm,right=2.5cm,top=2.5cm,bottom=2.5cm}
\usepackage{amsmath,amssymb,amsthm}
\usepackage{hyperref}
\usepackage{booktabs}
\usepackage{enumitem}
\usepackage{xltabular}
\hypersetup{colorlinks=true,linkcolor=blue,urlcolor=blue}

\title{方法论缺口识别与解决路径归档}
\author{李龙胜}
\date{\today}

\begin{document}

\maketitle

\begin{center}
\textit{
四卷体系重构完成后，发现四个方法论层面的薄弱点。\\
逐一分析后确认：两个已有理论工具覆盖，两个属于工程参数标定问题。\\
本文归档识别过程、解决路径和剩余工程任务，\\
作为维护工具包的补充文档。
}
\end{center}

\vspace{5mm}

\begin{abstract}
本文记录了四卷体系重构后识别出的四个方法论缺口：
从概念澄清到可计算模型的转化、贪心择优在策略互动中的最优性、
隐变量建模的方法论、以及验证与证伪的设计。
经过逐一分析，确认转化方法由六步闭环方法论覆盖，
贪心最优性已由贪心择优定理的非结合推广提供验证框架，
隐变量建模应切换回生成论的原生方法而非统计推断，
证伪框架内置于诚实标注的三级体系。
两个看似悬置的参数标定问题
（阈值公式中的 $\kappa_{\text{assoc}}, \Delta_{\min}, n_{\min}$ 等）
实为生产环境中自然涌现的实例化参数，非理论缺口。
\end{abstract}

\tableofcontents

\section{问题背景}

在完成四卷体系重构后，重新审视整个理论框架，
识别出四个方法论层面的薄弱点：
\begin{enumerate}
    \item \textbf{从概念澄清到可计算模型的转化方法}：
          强制去术语化擅长澄清概念骨架，但从骨架到可计算模型这一步的转化缺少系统化操作。
    \item \textbf{贪心择优在策略互动中的最优性}：
          贪心最优性定理在单方优化中已证明，
          但攻防双方同时贪心择优先的交互是否仍保证各自最优，尚未在博弈场景中严格验证。
    \item \textbf{隐变量建模的方法论}：
          $B$（攻击手法变化率）等隐变量应如何在生成论框架中建模？
          此前采用了统计推断路径，与生成论原生方法不一致。
    \item \textbf{验证与证伪的设计}：
          理论体系标注了层级，但缺乏系统的方法来设计验证或证伪实验，
          尤其当实验成本极高时。
\end{enumerate}

\section{缺失一：从概念澄清到可计算模型}

\textbf{识别}：
强制去术语化解决了“看清骨架”的问题，
但在安全运营场景中，从骨架到可计算的估计公式（如 $B$ 的量化映射）
缺乏系统化的操作步骤。

\textbf{解决路径}：
生成论非结合数学工具的方法论总结已提炼出六步闭环：
\begin{enumerate}
    \item 问题统一：将 B 不可直接观测识别为“单算符简并”问题。
    \item 工具定位：锁定非结合复合算符为打破简并的手段。
    \item 代数筛选：利用结合子选择定则筛选有效复合方向。
    \item 择优裁定：贪心择优选择成本最低的复合方案，直接确定 B 的取值。
    \item 最优保障：贪心择优定理的非结合推广提供最优性证明。
    \item 边界划定：在结合极限下退化为经典贪心定理。
\end{enumerate}
该六步闭环将概念澄清的结果直接输入第一步，
后续各步提供从代数结构到可计算模型的完整操作。

\textbf{状态}：已解决。方法是“生成论非结合数学工具的方法论总结”中的六步闭环。

\section{缺失二：贪心择优在策略互动中的最优性}

\textbf{识别}：
防御方贪心择优给出最优 $r^*$（假设攻击方固定），
攻击方贪心择优选择最优攻击尝试（假设防御方固定）。
双方同时贪心的交互是否等价于全局最优响应，尚未严格证明。

\textbf{解决路径}：
贪心择优定理的非结合推广已严格证明：
在结合子数量不超过阈值时，贪心等价于全局最优。
该阈值条件可以推广到攻防联合系统——
将双方操作视为联合生成步，验证联合系统的结合子数量是否在阈值内。
结合子选择定则的稀疏性使得该条件大概率自动满足。

\textbf{剩余工作}：
在安全运营场景中，将攻击方的告警生成和防御方的分析裁定
映射为八元数结构上的非结合复合操作，
计算联合结合子数量，与阈值比较。
这不是理论缺口，而是需要具体场景数据的形式化验证。

\textbf{状态}：理论框架已提供验证工具，实际验证需生产环境数据。

\section{缺失三：隐变量建模的方法论}

\textbf{识别}：
此前对 $B$ 的建模采用了基于 $\Delta D, \Delta N, \Delta \tau$ 的统计推断路径，
这与生成论处理挠率耦合 $\Lambda_T$ 和泄漏系数 $\varepsilon$ 的原生方法不一致。

\textbf{解决路径}：
生成论的原生方法不是从可观测数据中反推隐变量，
而是将隐变量的变化对应为生成步的代数结构变化，
通过遍历可能的复合算符、用贪心择优筛选最优方案，
直接从代数结构中提取隐变量取值。
$B$ 应切换到此方法：
将攻击手法的变异对应为告警生成步在八元数左理想上的复合变化，
通过结合子选择定则筛选有效变异方向，
贪心择优直接确定 $B$ 的最优取值。

\textbf{状态}：方向已切换。具体操作需在安全运营场景中建立
告警类型的八元数表示，并实现复合算符的遍历和筛选。

\section{缺失四：验证与证伪的设计}

\textbf{识别}：
诚实标注体系标注了各命题的层级（I/II/III），
但缺乏系统的方法来设计验证实验。

\textbf{解决路径}：
诚实标注本身内嵌了证伪框架：
\begin{itemize}
    \item \textbf{III级（启发式猜想）}：最易证伪——
          只需在红蓝对抗中失败一次，即需修正。
    \item \textbf{II级（定性正确，定量待标定）}：
          需要在实际场景中找到定性结论不成立的案例才能否证。
    \item \textbf{I级（严格证明）}：
          只能被逻辑错误或公设推翻所证伪。
\end{itemize}
该框架为所有命题提供了天然的证伪优先级排序，
不需要额外的“验证方法论”——
诚实标注本身就是在告诉审查者“应瞄准哪里”。

\textbf{状态}：已内置于诚实标注体系，无需额外方法论。

\section{工程参数标定（非方法论缺口）}

在分析缺失二和缺失三的过程中，
识别出若干参数（如 $\kappa_{\text{assoc}}, \Delta_{\min}, n_{\min}$）
需要在生产环境中标定。
这些是\textbf{实例化参数}，与从八元数乘法表中严格算出的
结构化参数（如 $\kappa = 4/7, \xi = 4/5$）有本质区别。
实例化参数不由生成论公设决定，
而是在具体生产过程中自动涌现的值。

\textbf{需要标定的关键参数}：
\begin{itemize}
    \item $\kappa_{\text{assoc}}$：安全运营场景中“每次结合子评估的记录成本”——
          对应分析师查台账的时间或AI做关联推理的算力消耗。
    \item $\Delta_{\min}$：候选复合方案间的最小差异——
          对应不同告警分类方案之间分析结果的分歧程度。
    \item $n_{\min}$：候选集中最小结合子数量——
          对应最简单的告警分析路径所需的关联步骤数。
    \item $\beta, \gamma$：同源关联优先定理中的上下文积累凸性参数
          和关联信息增益系数。
\end{itemize}

这些参数的值不由理论决定，而在生产部署后自然浮现。

\section{结论与后续工作}

四个方法论缺失已全部澄清。
生成论内部已有工具覆盖了概念转化、贪心推广、隐变量建模和验证框架。
剩余的工程任务是：
\begin{itemize}
    \item 在安全运营场景中建立告警类型的八元数表示，
          将攻击方生成步和防御方分析步映射为非结合复合操作。
    \item 验证联合系统的结合子数量是否在阈值内。
    \item 在生产部署中自动采集和标定实例化参数。
\end{itemize}

\vspace{5mm}
\noindent\textbf{文档信息}：
\begin{itemize}
    \item 作者：李龙胜
    \item 版本：第一版（方法论缺口归档）
    \item 许可：CC BY 4.0
    \item 前置依赖：
          四卷体系重构文档；
          贪心择优定理的非结合推广；
          生成论非结合数学工具的方法论总结
\end{itemize}

\end{document}